\documentclass[11pt]{article}
\usepackage[margin=25mm]{geometry}
\usepackage{fontspec,unicode-math}
\setmainfont{Latin Modern Roman}
\setmathfont{Latin Modern Math}
\usepackage{amsmath,amssymb,mathtools}
\usepackage{xurl,hyperref}
\hypersetup{colorlinks=true,urlcolor=blue,linkcolor=blue}
\setcounter{secnumdepth}{0}
\setlength{\parindent}{0pt}
\setlength{\parskip}{5pt}
\setlength{\emergencystretch}{3em}
\newcommand{\sref}[2]{\hyperlink{source:#1:#2}{[S#2]}}
\newcommand{\cataloguescope}[1]{\paragraph{Source question.}#1}
\newcommand{\floor}[1]{\left\lfloor#1\right\rfloor}
\begin{document}
\section{Convergence and Effective Rates of Unwinding Series}\label{statement}

\subsection{Definitions and mathematical statement}

\paragraph{Shared notation.}Standard number systems are denoted by $\mathbb N,\mathbb Z,\mathbb Q,\mathbb R,\mathbb C$; index ranges, computational models and specialized notation follow the cited problem source and the definitions below.


For a holomorphic function $F$ on the unit disk, write $F=B_0G_0$, extracting a Blaschke factor $B_0$ for its zeros. Iteratively factor $G_j-G_j(0)=B_{j+1}G_{j+1}$. The resulting unwinding expansion is formally $F=\sum_{j\geq0}G_j(0)B_0\cdots B_j$, with remainder obtained from the unfixed final $G_j$. Blaschke factors are inner factors associated with zeros in the disk.
\paragraph{Statement recorded in the supplied source.}
What can be rigorously established about convergence of the unwinding series, and can one obtain effective convergence rates for sufficiently regular functions $F$? \sref{statement}{2}
\paragraph{Scope qualification.}
The regularity class, convergence norm, and zero-factor conventions must be specified in a theorem. They are deliberately not fixed by the supplied research question.
\paragraph{Source status.}
The supplied catalogue labels this record partially solved. Its source-status metadata was checked against the supplied snapshot on 2026-10-11. This import does not independently establish current open status or mathematical correctness. The catalogue literature assessment is dated 2026-08-22.
\subsection{Short English statement}
Prove convergence and quantitative approximation rates for the iterative zero-factor expansion.
\subsection{Sources}
\begin{itemize}
\item[\hypertarget{source:statement:1}{S1}] Ulam AI and the UnsolvedMath Contributors, UnsolvedMath: A Curated Collection of Open Mathematics Problems, catalogue entry. Catalogue curation and metadata: CC BY 4.0. The supplied source question is retained; definitions and the English summary are editorial additions.
\url{https://huggingface.co/datasets/ulamai/UnsolvedMath}.
\item[\hypertarget{source:statement:2}{S2}] Applied Harmonic Analysis and Data Science (2022), Oberwolfach Reports, DOI 10.4171/owr/2021/55, pp. 3007--3066
\url{https://doi.org/10.4171/owr/2021/55}.
\item[\hypertarget{source:statement:3}{S3}] Applied Harmonic Analysis and Data Science, Oberwolfach Reports 55/2021, primary report PDF
\url{https://ems.press/content/serial-article-files/46933}.
\textit{Underlying Oberwolfach Reports material retains its source-specific licence, including CC BY-SA 4.0 where applicable; see THIRD\_PARTY\_NOTICES.md.}
\end{itemize}
\paragraph{Preparation and review.}
Original question attributed to the cited source and UnsolvedMath contributors. Codex (GPT-6; exact serving revision not exposed) assisted with definitions, English summaries, source inspection and formatting. No new human mathematical authorship or independent proof verification is asserted. Catalogue identity checked against the complete supplied registry and existing repository statements; known flagged repetitions and obvious identical conjectures excluded. Full mathematical-equivalence/literature review has not been performed. Statement preparation and final review followed contribution prompts C and D; automated checks verify structure and identity, not open status.
\label{end:statement}
\end{document}
